\section*{Chapter \ref{ch:iv1:timed-test-1} --- Timed Test 1}
\markboth{\protect\omTocMark{}{Answers: Timed Test 1}}{\protect\omTocMark{}{Answers: Timed Test 1}}

\begin{solution}{iq:iv1:timed-test-1:1}
The answer is $144/1024 = 9/64 \approx 0.141$.

Let $a_n$ be the number of strings of $n$ flips with no two adjacent heads. A good string ends either in tails, after any good string of length $n - 1$, or in tails-heads, after any good string of length $n - 2$. So $a_n = a_{n-1} + a_{n-2}$\index{Fibonacci recurrence}, with $a_1 = 2$ and $a_2 = 3$. The sequence runs $2, 3, 5, 8, 13, 21, 34, 55, 89, 144$, and there are $2^{10} = 1024$ equally likely strings.

\iqwhy{Writing down the last flip or two is the fastest route through most counts of this kind, and it avoids a long case analysis.}
\end{solution}

\begin{solution}{iq:iv1:timed-test-1:2}
The first player wins with probability $4/9$.

Use first-step analysis\index{first-step analysis}, with two states for the player about to flip. Let $y$ be that player's chance of winning when the flip before was heads, and $x$ when it was not. If the flip before was heads, tails wins at once, and heads hands the same state to the opponent:
\[
y = \tfrac12 + \tfrac12(1 - y), \quad\text{so } y = \tfrac23.
\]
If not, any flip hands the turn over, with heads giving the opponent the good state:
\[
x = \tfrac12(1 - y) + \tfrac12(1 - x), \quad\text{so } x = \tfrac49.
\]
The first player starts in the state $x$.

\iqwrong{Moving first is not an advantage here. Your first flip cannot win, and it may hand the opponent a winning chance.}
\end{solution}

\begin{solution}{iq:iv1:timed-test-1:3}
The answer is 14 drops, starting from floor 14.

If your plan allows $k$ drops, the first drop can be at most from floor $k$: if the egg breaks, the second egg has to walk up through the $k - 1$ floors below, one at a time. If it survives, $k - 1$ drops remain, so the next jump is $k - 1$ floors, and so on. Thus $k$ drops can cover $k + (k - 1) + \dots + 1 = k(k + 1)/2$ floors\index{triangular numbers}. Since $13 \cdot 14/2 = 91$ is too few and $14 \cdot 15/2 = 105$ is enough, $k = 14$, with the first egg going to floors 14, 27, 39, and so on.

\iqwrong{Equal steps of 10, which costs up to 19 drops. Equal steps waste the cost of the early drops.}
\iqfollow{And with 3 eggs?}{9 drops.}
\end{solution}

\begin{solution}{iq:iv1:timed-test-1:4}
The answer is $5/16$.

Find it by areas in the unit square\index{geometric probability}. The event $|X - Y| > a$ is two right-angled corner triangles, each with legs $1 - a$, with total area $(1 - a)^2$. So $P(|X - Y| < a) = 1 - (1 - a)^2$, and
\[
P\bigl(\tfrac14 < |X - Y| < \tfrac12\bigr) = \bigl(\tfrac34\bigr)^2 - \bigl(\tfrac12\bigr)^2 = \tfrac{9}{16} - \tfrac{4}{16} = \tfrac{5}{16} ,
\]
which is the shaded area in \cref{fig:iv1:band}.

\begin{omfigure}
\begin{tikzpicture}[x=3.6cm,y=3.6cm]
\fill[omDef!25] (0,0.25) -- (0,0.5) -- (0.5,1) -- (0.75,1) -- cycle;
\fill[omDef!25] (0.25,0) -- (0.5,0) -- (1,0.5) -- (1,0.75) -- cycle;
\draw (0,0) rectangle (1,1);
\draw[gray,dashed] (0,0) -- (1,1);
\draw[omDef] (0,0.25) -- (0.75,1);
\draw[omDef] (0,0.5) -- (0.5,1);
\draw[omDef] (0.25,0) -- (1,0.75);
\draw[omDef] (0.5,0) -- (1,0.5);
\node[left,font=\footnotesize] at (0,0.25) {$\tfrac14$};
\node[left,font=\footnotesize] at (0,0.5) {$\tfrac12$};
\node[left,font=\footnotesize] at (0,1) {$1$};
\node[below,font=\footnotesize] at (0.25,0) {$\tfrac14$};
\node[below,font=\footnotesize] at (0.5,0) {$\tfrac12$};
\node[below,font=\footnotesize] at (1,0) {$1$};
\node[below left,font=\footnotesize] at (0,0) {$0$};
\node[right,font=\footnotesize] at (1.02,0.5) {$X$ across, $Y$ up};
\end{tikzpicture}

\omcaption{The unit square with the two shaded bands where $|X - Y|$ lies between $\tfrac14$ and $\tfrac12$. The dashed line is $X = Y$.}\label{fig:iv1:band}
\end{omfigure}
\end{solution}

\begin{solution}{iq:iv1:timed-test-1:5}
About $0.39$ of months lose money, or $0.386$ to three places. Anything from $0.38$ to $0.39$ is fine by hand.

Over 12 independent months, mean and variance both scale by 12, so the annual Sharpe ratio is $\sqrt{12}$ times the monthly one\index{Sharpe ratio}. The monthly Sharpe ratio is therefore $1/\sqrt{12} \approx 0.29$. A month loses money when a standard normal falls below $-0.29$. Near zero the normal density is about $0.4$, so this probability is about $0.5 - 0.29 \cdot 0.4 = 0.384$, and the exact value is $0.386$.

\iqwhy{Even a strategy with a Sharpe ratio of 1 loses in almost two months out of five. Expect that before you judge a live track record by a bad quarter.}
\end{solution}

\begin{solution}{iq:iv1:timed-test-1:6}
The determinant is $81/128 \approx 0.633$.

Write the matrix as $0.75\,I + 0.25\,\mathbf{1}\mathbf{1}^\top$, where $\mathbf{1}$ is the vector of ones\index{eigenvalues of a low-rank update}. The matrix $\mathbf{1}\mathbf{1}^\top$ has eigenvalue 5 in the direction of $\mathbf{1}$ and eigenvalue 0 on the four directions orthogonal to it. So the eigenvalues are $0.75 + 0.25 \cdot 5 = 2$, once, and $0.75$, four times. The determinant is the product: $2 \cdot (3/4)^4 = 81/128$.

\iqfollow{How negative could the common off-diagonal entry $c$ be, for this to remain a correlation matrix?}{The large eigenvalue becomes $1 + 4c$, which must not be negative. So $c \ge -1/4$.}
\end{solution}

\begin{solution}{iq:iv1:timed-test-1:7}
The probability is $1/99$. All totals from 1 to 99 are equally likely.

This is a P\'olya urn\index{P\'olya urn}: each outcome adds one to its own count, and the next probability is the proportion. Prove the claim by induction. After $n$ shots, suppose the number of makes is uniform on $1, \dots, n - 1$, so each value has probability $1/(n - 1)$. To make $k$ of $n + 1$, either you had $k - 1$ and made the shot, or you had $k$ and missed:
\[
\frac{1}{n - 1}\left(\frac{k - 1}{n} + \frac{n - k}{n}\right) = \frac{1}{n - 1}\cdot\frac{n - 1}{n} = \frac1n .
\]
So after $n + 1$ shots the number of makes is uniform on $1, \dots, n$. It starts true at $n = 2$, and with 100 shots there are 99 values.

\iqwrong{Expecting a peak at 50 like a binomial. The shots are not independent, and an early streak persists.}
\end{solution}

\enlargethispage{-4\baselineskip}%
\begin{solution}{iq:iv1:timed-test-1:8}
The expected longest piece is $11/18 \approx 0.611$.

Use the tail formula\index{tail-sum formula} with inclusion--exclusion. Any one piece exceeds $x$ with probability $(1 - x)^2$. Two chosen pieces both exceed $x$ with probability $(1 - 2x)^2$ when $2x < 1$, and all three with probability $(1 - 3x)^2$ when $3x < 1$. Hence
\[
P(L > x) = 3(1 - x)^2 - 3(1 - 2x)_+^2 + (1 - 3x)_+^2 .
\]
Integrating over $x$ from 0 to 1 gives $1 - \tfrac12 + \tfrac19 = \tfrac{11}{18}$.

\iqfollow{Is there a quicker route?}{The smallest piece has mean $\tfrac13 \cdot \tfrac13$, the next is larger by $\tfrac13 \cdot \tfrac12$ on average, and the largest by a further $\tfrac13 \cdot 1$. So the largest has mean $\tfrac13(1 + \tfrac12 + \tfrac13) = \tfrac{11}{18}$. With $n$ pieces it is $H_n/n$, where $H_n$ is the harmonic sum.}
\end{solution}

\begin{solution}{iq:iv1:timed-test-1:9}
The answer is 7 strips.

Each strip has three possible stories: positive in round 1, positive in round 2, or never positive. With $s$ strips there are at most $3^s$ outcomes, and each bottle needs its own outcome\index{information bound}. Since $3^6 = 729 < 1000 \le 2187 = 3^7$, six strips are not enough, even with choices in round 2 that depend on round 1: if $k$ strips test positive in round 1, only $2^{s-k}$ outcomes remain, and $\sum_k \binom{s}{k}2^{s-k} = 3^s$ all the same.

Seven are enough. Write each bottle's number in base 3 with 7 digits. Strip $i$ tests a bottle in round 1 if digit $i$ is 0, in round 2 if it is 1, and not at all if it is 2. The strips that fire, and when, read out the poisoned bottle's digits.

\iqwrong{The answer 10, from $2^{10} \ge 1000$, which ignores the second round entirely.}
\end{solution}

\newpage
\begin{solution}{iq:iv1:timed-test-1:10}
You get wet on $21/370 \approx 5.68\%$ of trips.

Model the number $i$ of umbrellas where you are\index{stationary distribution}, with $r = 3$ in total, rain probability $p = 0.3$, and $q = 0.7$. With none, the next place holds all $r$. With $i > 0$, the next place holds $r - i$ if it stays dry, and $r - i + 1$ if it rains and you carry one over. Solving the balance equations gives
\[
\pi_0 = \frac{q}{r + q}, \qquad \pi_i = \frac{1}{r + q} \quad (i = 1, \dots, r).
\]
You are wet when you have none and it rains, so the fraction is $p\,\pi_0 = 0.3 \cdot 0.7/3.7 = 21/370$.

\iqwrong{Using a uniform distribution, $\pi_0 = 1/4$, which gives $0.075$. State zero is rarer than the others, by the factor $q$.}
\end{solution}

\begin{solution}{iq:iv1:timed-test-1:11}
A survives with probability $25/63 \approx 0.397$, by firing into the air on the first turn.

Work backwards. In a duel between A and B with A firing first, a full round passes with no hit with probability $\tfrac23 \cdot \tfrac13 = \tfrac29$, so A wins with probability $\tfrac13/(1 - \tfrac29) = \tfrac37$.

With all three alive, B must aim at C, since otherwise C kills B at once. C, if reached, aims at B, the better shot. So if A fires in the air, then B hits C with probability $2/3$, leaving A a duel with the first shot, worth $3/7$. Otherwise C kills B, and A has one shot at C, worth $1/3$. Altogether
\[
\tfrac23 \cdot \tfrac37 + \tfrac13 \cdot \tfrac13 = \tfrac{25}{63}.
\]
Aiming at C is worse: a hit leaves A to take B's fire first, worth only $1/7$, for a total of $59/189$. Aiming at B is worst, for a total of $50/189$, since hitting B means C kills A\index{backward induction}.

\iqwhy{The weakest player does best by not shooting, because the other two have to deal with each other first.}
\end{solution}

\begin{solution}{iq:iv1:timed-test-1:12}
You need about 194 flips. Anything from 190 to 200 is fine.

Let $n$ be the number of flips, and work with the proportion of heads\index{power of a test}. The test rejects when the proportion is further than $1.96 \cdot 0.5/\sqrt n$ from $0.5$. Under the alternative, the proportion has mean $0.6$ and standard deviation $\sqrt{0.24}/\sqrt n \approx 0.49/\sqrt n$. For 80\% power, the rejection point must lie $0.84$ of those standard deviations below $0.6$. So
\[
0.1 = \frac{1.96 \cdot 0.5 + 0.84 \cdot 0.49}{\sqrt n} = \frac{1.39}{\sqrt n}, \qquad n \approx 193 .
\]
Keeping more digits gives 194, and the exact binomial power at 194 flips is $0.81$. The usual shortcut of $0.5$ for both standard deviations gives 196, which is fine.

\iqwrong{Using $1.645$, the one-sided value, gives about 152. That would be the answer to another question.}
\end{solution}
